%% Trame série asynchrone (UART) d'un caractère : repos à 1, bit de start à 0,
%% 8 bits de donnée poids faible (D0) en tête, bit de parité, bit de stop, retour au repos.
%% Motif tracé : D0..D7 = 1 0 1 0 1 1 0 0 (donnée 0011 0101 lue du msb au lsb).
%% Largeur d'un intervalle : 0,58 cm.
\documentclass[tikz,border=4pt]{standalone}
\usepackage{liaisons}
\begin{document}
\begin{tikzpicture}[x=1cm,y=1cm,
  sep/.style={black!40, dash pattern=on 1.6pt off 1.6pt, line width=0.5pt},
  grad/.style={font=\scriptsize},
  champ/.style={solideD, line width=0.9pt, line join=round}]

%% ================== niveaux logiques ==================================
\node[grad, anchor=east, inner sep=3pt] at (-0.55,0)    {0};
\node[grad, anchor=east, inner sep=3pt] at (-0.55,0.75) {1};

%% ================== séparateurs d'intervalles =========================
\foreach \i in {0,1,...,11} { \draw[sep] ({\i*0.58},-0.30) -- ({\i*0.58},0.95); }

%% ================== tracé de la trame =================================
\draw[solideB, line width=1.4pt, line join=round]
  (-0.50,0.75) -- (0,0.75) -- (0,0) -- (0.58,0) -- (0.58,0.75) -- (1.16,0.75) --
  (1.16,0) -- (1.74,0) -- (1.74,0.75) -- (2.32,0.75) -- (2.32,0) -- (2.90,0) --
  (2.90,0.75) -- (4.06,0.75) -- (4.06,0) -- (5.22,0) -- (5.22,0.75) -- (6.88,0.75);

%% ================== états de repos ====================================
\node[grad, text=black!60, anchor=south] at (-0.25,0.80) {Repos};
\node[grad, text=black!60, anchor=south] at (6.63,0.80) {Repos};

%% ================== noms des bits de donnée ===========================
\foreach \k in {0,1,...,7} {
  \node[grad, anchor=north, inner sep=2pt] at ({0.58*\k+0.87},-0.30) {D\k}; }

%% ================== champs de la trame ================================
%% chaque champ : un crochet ouvert vers le haut, sous l'axe des bits
\foreach \a/\b in {0/0.58, 0.58/5.22, 5.22/5.80, 5.80/6.38} {
  \draw[champ] (\a,-0.75) -- (\a,-0.85) -- (\b,-0.85) -- (\b,-0.75); }
\node[grad, text=solideD, anchor=north, inner sep=3pt] at (0.29,-0.88) {start};
\node[grad, text=solideD, anchor=north, inner sep=3pt] at (2.90,-0.88)
      {8 bits de donnée (D0 en tête)};
\node[grad, text=solideD, anchor=north, inner sep=3pt] at (5.51,-0.88) {parité};
\draw[champ] (6.09,-0.85) -- (6.09,-1.22);
\node[grad, text=solideD, anchor=north, inner sep=1pt] at (6.09,-1.22) {stop};
\end{tikzpicture}
\end{document}
